% \section{Linearização e Discretização}
% Considerando o modelo dinâmico contínuo do sistema manipulador acoplado à base, sem contatos ativos, dado por:
% \begin{equation}
% \mathbf{M}(\vec{q})\ddot{\vec{q}}
% +\mathbf{C}(\vec{q}, \dot{\vec{q}})\dot{\vec{q}}
% +\vec{g}(\vec{q})
% =
% \vec{\tau}
% -\mathbf{D} \dot{\vec{q}}
% +\mathbf{J}(\vec{q})^T \vec{\lambda}
% \label{eq:modelo_dinamico}
% \end{equation}
% Define-se o vetor de estado como $\vec{x} = \begin{bmatrix} \vec{q} \\ \dot{\vec{q}} \end{bmatrix}$ e a entrada de controle como $\vec{u} = \vec{\tau}$. Um ponto de operação $(\vec{x}^*, \vec{u}^*)$ satisfaz a condição de equilíbrio:
% \begin{equation}
% \dot{\vec{x}}^* = \vec{0}, \quad \mathbf{C}(\vec{q}^*, \vec{0}) + \vec{g}(\vec{q}^*) = \vec{\tau}^* + \mathbf{J}(\vec{q}^*)^T \vec{\lambda}^*
% \label{eq:equilibrio}
% \end{equation}
% Para manobras sem contatos, pode-se assumir $\vec{\lambda}^* = \vec{0}$. A forma explícita do modelo é obtida isolando-se $\ddot{\vec{q}}$:
% \begin{equation}
% \ddot{\vec{q}} = \mathbf{M}^{-1}(\vec{q})\left( \vec{u} - \mathbf{D} \dot{\vec{q}} - \mathbf{C}(\vec{q}, \dot{\vec{q}})\dot{\vec{q}} - \vec{g}(\vec{q}) \right)
% \label{eq:forma_explicita}
% \end{equation}
% Substituindo em um sistema de primeira ordem, obtém-se:
% \begin{equation}
% \dot{\vec{x}} = \begin{bmatrix} \dot{\vec{q}} \\ \vec{f}(\vec{q}, \dot{\vec{q}}, \vec{u}) \end{bmatrix} = \vec{F}(\vec{x}, \vec{u})
% \label{eq:sistema_estado}
% \end{equation}
% Caso haja vínculos diferenciais holonômicos, como a normalização de quaternions unitários (por exemplo, $\|\vec{q}_b\|^2 = 1$), impõe-se a condição:
% \begin{equation}
% \vec{\Phi}(\vec{q}) = \vec{0}
% \label{eq:vinculo_phi}
% \end{equation}
% A linearização dessas restrições no ponto $\vec{q}^*$ resulta em:
% \begin{equation}
% \delta \vec{\Phi} = \mathbf{J}_{\Phi}(\vec{q}^*) \, \delta \vec{q} = \vec{0}.
% \label{eq:linearizacao_vinculo}
% \end{equation}
% Com $\mathbf{J}_{\Phi} = \partial \vec{\Phi} / \partial \vec{q}$. Pode-se então projetar a dinâmica linearizada nesse subespaço nulo, eliminando os multiplicadores de Lagrange, ou manter o sistema como uma DAE e aplicar regularização (ex: método de Baumgarte).
% Para implementação digital, adota-se discretização por zero-order hold (ZOH) com período de amostragem $T_s$. A forma contínua linearizada:
% \begin{equation}
% \dot{\vec{x}} = \mathbf{A} \vec{x} + \mathbf{B} \vec{u}.
% \label{eq:sistema_linear_continuo}
% \end{equation}
% É convertida em:
% \begin{equation}
% \vec{x}_{k+1} = \mathbf{A}_d \vec{x}_k + \mathbf{B}_d \vec{u}_k.
% \label{eq:sistema_discreto}
% \end{equation}
% Com as transformações:
% \begin{equation}
% \mathbf{A}_d = e^{\mathbf{A} T_s}, \quad \mathbf{B}_d = \int_0^{T_s} e^{\mathbf{A} \tau} d\tau \, \mathbf{B}
% \label{eq:discretizacao_AdBd}
% \end{equation}
% Como critério prático, recomenda-se que $T_s$ seja no máximo $\frac{1}{10}$ do tempo dominante do sistema. De modo equivalente, utiliza-se:
% \begin{equation}
% T_s \leq \frac{\pi}{10 \, \omega_{\text{bw}}}.
% \label{eq:criterio_Ts}
% \end{equation}
% Onde $\omega_{\text{bw}}$ é a frequência de banda do sistema. Em sistemas SISO auxiliares (ex: filtros), pode-se adotar o método bilinear de Tustin.
% Por fim, os pares $(\mathbf{A}, \mathbf{B})$ devem ser estimados numericamente por diferenças finitas no ponto de operação. A discretização deve ser exata (ZOH) para garantir precisão. Em sistemas com quaternions, é necessário impor normalização periódica ou aplicar projeção a cada passo discreto, garantindo estabilidade numérica e coerência física.






%==================================================
%----- Linearização em tempo continuo
%==================================================
\section{Linearização em tempo contínuo}
O modelo contínuo com vínculos holonômicos, conforme apresentado em
\eqref{sec:restricoes_holonomicas},
e formulado com os multiplicadores de Lagrange foi apresentado em
\eqref{eq:dinamica_mult_lagrange},
juntamente com as restrições escritas de
\eqref{eq:PhiJ1} a \eqref{eq:PhiJ2},
e seus níveis diferencial e de aceleração em
\eqref{eq:vellevel} a \eqref{eq:acclevel}.
Quando há dissipação de Rayleigh, utiliza-se a forma aumentada presente em
\eqref{eq:dae_diss},
enquanto que o bloco correspondente é
\eqref{eq:block_diss}.
Nesta seção, adota-se a notação de estado para obter o modelo linearizado a ser utilizado nas leis de controle e análise local.

\subsection{Definição de estado e ponto de operação}
Considerando o vetor de estado sendo:
\begin{equation}
\vec x=
\begin{bmatrix}
\vec q\\
\dot{\vec q},
\end{bmatrix}
\end{equation}
e a entrada sendo $\vec u=\vec\tau$;
como $\dot{\vec q}^{}=\ddot{\vec q}^{}=\vec 0$, resulta 
$\vec u^{}+\mathbf J(\vec q^{},t)^{\top}\vec\lambda^{}=\vec g(\vec q^{})$
\fn{Ver \eqref{eq:dinamica_mult_lagrange}};
sem contato, considera-se $\vec\lambda^{}=\vec 0$ e $\vec u^{}=\vec g(\vec q^{*})$.

\subsection{Forma de primeira ordem}
A passagem para primeira ordem é a mesma empregada na montagem do bloco de acelerações em
\eqref{eq:sist_linear_blocos},
isola-se $\ddot{\vec q}$ de
\eqref{eq:dinamica_mult_lagrange}, mas
caso haja dissipação, $\mathbf D$,
e ejusto com $\dot{\vec q}$,
obtém-se $\dot{\vec x}=\vec F(\vec x,\vec u,\vec\lambda)$.
Como as restrições são tratadas via estabilização de Baumgarte, a contribuição em nível de aceleração é dada por \eqref{eq:acclevel_baumgarte};
caso contrário, o acoplamento com $\vec\lambda$ permanece como em \eqref{eq:sist_linear_blocos}.

\subsection{Linearização}
A variação em torno de $(\vec x^{*},\vec u^{*})$ é:
\begin{equation}
\delta\dot{\vec x}
=
\mathbf A\,\delta\vec x+\mathbf B\,\delta\vec u,
\label{eq:ss_linear_cont1}
\end{equation}
sendo:
\begin{equation}
\mathbf A=\left.\frac{\partial \vec F}{\partial \vec x}\right|_{*},
\label{eq:ss_linear_cont2}
\end{equation}

e

\begin{equation}
\mathbf B=\left.\frac{\partial \vec F}{\partial \vec u}\right|_{*}.
\label{eq:ss_linear_cont}
\end{equation}

Pela estrutura mecânica, escreve-se:
\begin{equation}
\mathbf A=
\begin{bmatrix}
\mathbf 0 & \mathbf I\\
\mathbf A_q & \mathbf A_v
\end{bmatrix},
\qquad
\mathbf B=
\begin{bmatrix}
\mathbf 0\\
\mathbf M(\vec q^{*})^{-1}
\end{bmatrix},
\label{eq:blocos_A_B}
\end{equation}
onde
$\mathbf A_q=\left.\dfrac{\partial \vec f}{\partial \vec q}\right|_{*}$ e
$\mathbf A_v=\left.\dfrac{\partial \vec f}{\partial \dot{\vec q}}\right|_{*}$.
Quando em equilíbrio, $\dot{\vec q}^{*}=\vec 0$, os termos de Coriolis/centrífuga que são proporcionais a
$\dot{\vec q}$ e se anulam, resultando em uma aproximação local:
\begin{equation}
\mathbf A_q \approx
-\,\mathbf M(\vec q^{*})^{-1}
\left.\frac{\partial \vec g}{\partial \vec q}\right|_{\vec q^{*}},
\qquad
\mathbf A_v \approx
-\mathbf M(\vec q^{*})^{-1}\mathbf D,
\label{eq:aprox_Aq_Av}
\end{equation}
sendo $\mathbf D=\mathbf 0$
quando não se emprega dissipação de Rayleigh, vide \eqref{eq:Rayleigh} até \eqref{eq:Qd}.

\subsection{Tratamento dos vínculos na linearização}
Para $\Phi(\vec q)=\vec 0$, conforme \eqref{eq:PhiJ1},
a linearização do vínculo é
$\delta\Phi=\mathbf J_{\!\Phi}(\vec q^{*})\,\delta\vec q=\vec 0$,
já estabelecida em \eqref{eq:vellevel} e
\eqref{eq:acclevel}.
Duas opções são consistentes com as seções anteriores:
(i) Projeção: escolha da base $\mathbf N$ tal que
$\mathbf J_{\!\Phi}(\vec q^{*})\,\mathbf N=\mathbf 0$ e
$\delta\vec q=\mathbf N\,\vec\eta$;
então $\delta\dot{\vec z}=\mathbf A_r\,\delta\vec z+\mathbf B_r\,\delta\vec u$
com $\delta\vec z=\begin{bmatrix}\vec\eta\\ \dot{\vec\eta}\end{bmatrix}$
e $\mathbf A_r,\mathbf B_r$
obtidas por mudança de base, sem derivar novamente
\eqref{eq:sist_linear_blocos}.  
(ii) Baumgarte: incorporar
\eqref{eq:acclevel_baumgarte} como relação de aceleração adicional, equivalente a um controle PD sobre $\Phi$ em \eqref{eq:PD_phi}, preservando a estrutura DAE.

\subsection{Observações}
A simulação e o controle serão em tempo contínuo, portanto, não será utilizada a forma discretizada.
% $\,\vec x_{k+1}=\mathbf A_d\vec x_k+\mathbf B_d\vec u_k$.
Para o cálculo prático de
$\mathbf A$ e $\mathbf B$,
podem-se empregar diferenças finitas ao redor de
$(\vec x^{*},\vec u^{*})$,
assegurando a normalização do quaternion e a projeção do par
$(\delta\vec q,\delta\dot{\vec q})$
conforme \eqref{eq:PhiJ1}.
